\subsection{The Quantum Computing Edge in Quantum Chemistry}
% Motivate electronic-structure problems as a natural target for quantum computers.
% Explain why classical scaling becomes difficult for correlated molecules and excited states.
% Position near-term algorithms as practical approximations rather than fault-tolerant solutions.

\subsection{Quantum Subspace Expansion for Excited States}
% Introduce QSE as a post-VQE subspace method for ground and excited states.
% Explain the idea of applying excitation/operator pools to a reference state.
% Connect QSE to projected Hamiltonian and overlap matrices.

\subsection{Numerical Instability in Practical QSE}
% State the central practical problem: the formal QSE space can contain weakly supported directions.
% Explain why small overlap eigenvalues make the generalized eigenvalue problem sensitive.
% Motivate stability filtering as necessary but potentially lossy.

\subsection{Generalized Eigenvalue Conditioning}
% Introduce the role of the overlap matrix spectrum in QSE stability.
% Preview thresholds, condition numbers, retained rank, and stable dimension.
% Emphasize the trade-off between numerical stability and subspace completeness.

\subsection{Current QSE Studies and Contributions of This Work}
% Summarize prior QSE work on excited states, response properties, and NISQ implementations.
% Identify the gap: practical stability and completeness are not always characterized together.
% State this work's focus on QSE instability, representability of the projected space, finite-shot effects, and scaling.
